\section{Stable sorting under a strict square saving}\label{sec:sorting}

A tensor network may have to move a table to the place where it will be
used, or bring together all the queries addressed to one table. Moving a table involves long records, while gathering its queries
requires sorting many short ones. We show
that a strict improvement on the exponent $9/4$ would make both operations
cheap enough for the final assembly in Sections~\ref{sec:assembly}
and~\ref{sec:shaving}.

Fix the finite field $k$ from Section~\ref{sec:obstruction} and let $n=2^h$.
The tagged permutation tensor from Section~\ref{sec:list} allows the
destination of every record to be part of the input. Write $j_i$ for
the $h$-bit destination of the $w$-bit source record $x_i$, and $y_j$
for the proposed record at position $j$. Then $\mathrm{Tag}(h,w)$ is the
indicator of the event that the $j_i$ are distinct and $y_{j_i}=x_i$
for every $i$. Similarly, the stable-sort tensor is the indicator that a
proposed output sorts the input records by a designated key and preserves
the input order within each key. Both indicators include the payload bits
in the equality tests.

\begin{theorem}[Conditional stable sorting]\label{thm:sort}
Suppose $\omega_k<9/4$, and let $n=2^h$.
\begin{enumerate}
\item For some $1<p<2$, $\mathrm{Tag}(h,h)$ has an exact trivalent
network over $k$ of size $O(nh^p)$. The same bound holds for stable
sorting of records of at most $ch$ bits, for every fixed $c$.
\item For every $\zeta>0$ there is $a>1$ such that, for every $w\ge h^a$,
both $\mathrm{Tag}(h,w)$ and stable sorting of $w$-bit records have
exact trivalent networks over $k$ of size $O_\zeta(nw h^\zeta)$.
\end{enumerate}
\end{theorem}

\subsection{Evaluating a table by sorting its queries}

At a station of arity $t$, the incident bits give the address of an
entry in a table of length $2^t$. When many copies of the station are
evaluated together, only the multiplicity of each address matters to
the product of the selected entries. Sorting the queries recovers
these multiplicities and spreads the cost of the table over the batch.

\begin{lemma}[Batched table evaluation]\label{lem:batch-table}
Assume the ordinary-width stable-sort bound of Theorem~\ref{thm:sort}
over $k$, with exponent $p>1$. Let a network of binary wires and
$K$-valued tables represent a tensor $T$, where $K/k$ is any field
extension. If $t_v\le u$ is the arity of station $v$ and
$s=\sum_v(1+t_v)$, then $T^{\otimes 2^u}$ has an exact trivalent
$K$-network of size
\[
                  O\bigl(2^u s(1+u)^{p-1}\bigr).
\]
The constant is independent of $K$ and of the entries in the tables.
\end{lemma}

\begin{proof}
First consider one table $f\colon\{0,1\}^t\to K$, with $t\ge1$, and
expose the $2^t$ query addresses $x_1,\ldots,x_{2^t}$ as tensor legs.
If $c_z$ is the number of occurrences of address $z$, the required
product is $\prod_z f(z)^{c_z}$, with an empty product at an unused
address. Create one dictionary row for each $z$ and guess its count in
$t+1$ bits. On count bit $j$, for $0\le j\le t$, place a unary factor
with entries
\[
                    (1,\ f(z)^{2^j}).
\]
The position of this factor already determines $z$. Thus its two entries
are fixed coefficients of the network.

Give each query its own row, with zeros in the count field. Sort the
$2^{t+1}$ rows by address and then by row type, so that the dictionary
row follows all queries to that address. A counter, initially zero,
increments on every query. On a dictionary row it must agree with the
carried count and then resets to zero. This checks the count even at
addresses with no queries. For fixed exposed addresses there is exactly
one valid choice of counts, sorted rows and counter states. The
contraction therefore equals $\prod_i f(x_i)$ in every characteristic,
also when a table entry vanishes.

The counter uses $O(2^t(1+t))$ gates by ripple-carry addition.
Sorting the $O(t+1)$-bit rows costs $O(2^t(1+t)^p)$. The sorting
network realizes a zero--one tensor over $k$ and hence realizes the
same tensor over $K$ by scalar extension.

Group the $2^u$ copies of each station of arity $t$ into $2^{u-t}$
batches of $2^t$. A station of arity zero supplies a scalar directly.
Connect the exposed wires according to the original network in each
copy. The result contracts to $T^{\otimes 2^u}$ and its size is
bounded by a constant times
\[
  2^u\sum_v(1+t_v)^p
       \ \le\ 2^u(1+u)^{p-1}\sum_v(1+t_v).
\]
\end{proof}


\subsection{The two scales of the sorting proof}\label{sec:sort-road}

Assume $\omega_k<9/4$ throughout Part~II. We first route wide records.
A hop that permutes records within a smaller address set can be handled
recursively. The objective is to shorten the address by an arbitrarily
large constant factor while keeping the total width transported in each summand within a
constant multiple of the input width, independently of that factor. Repeating such
a reduction gives the wide-record estimate in Theorem~\ref{thm:sort}.
Once that estimate is available, two levels of radix sorting give the
ordinary-width estimate. A short key selects a digit within a block,
and whole packs of records with the same digit then travel together.

Section~\ref{sec:sort-action} uses the incidence inequalities of Part~I
to compare the price of a target permutation with the prices of hops on
smaller sets of addresses. The comparison first holds for an action on
a CAT(0) space. Local price data sufficient to realize such an action
give the same bound. If the data fail a local consistency test, spectral
comparison instead supplies a tensor identity that reduces the cost.
Section~\ref{sec:sort-route} combines these two conclusions to route
payloads and to bound the number of recursive calls in a fixed network.

Throughout Parts~II--IV we use base-two logarithms for ordinary prices.
In these units the distance $d$ is the price of carrying one bit.
On a fixed finite list of sites, price data with bounded scale,
including scale zero, form a compact set by \eqref{eq:prices-range}
and Section~\ref{sec:prices-background}. A bound on diameter suffices
here. Equality along a spanning tree of binary wires realizes the
full diagonal and bounds its price.

Distances between sites obtained by finitely many geodesic interpolations
at fixed fractions depend continuously on the original prices. The
enlarged configurations still have bounded scale. Along any convergent
subnet the interpolation equalities pass to the limit, where
Lemma~\ref{lem:geodesic} determines the distances from the original data.


\section{From incidence to permutation actions}\label{sec:sort-action}

A permutation can be factored into hops that change only a short
part of the address. Adding their prices gives a bound proportional
to the number of factors. Sorting with wide records requires a
bound independent of that number. We first obtain such a bound on
displacement in a metric space. We then show that finite tests on
the addresses suffice to bound the mean price of the payload wires.

\subsection{Two incidence inequalities}

Fix a power $q\ge100$ of the characteristic of $k$. In a placement
of a bipartite incidence graph, a vertex on either side is assigned
a site. The incidence energy $e$ is the mean squared distance between
sites at the ends of a uniform incidence.

\begin{lemma}\label{lem:sort-inc}
Suppose $\omega_k<9/4$. In any complete $\operatorname{CAT}(0)$
spectral space over $k$, consider a placement of either of the
following graphs.
\begin{enumerate}
\item The points and lines of the projective space $\mathrm{PG}(\ell-1,q)$,
where $\ell\ge3$.
\item The points of $\mathbb F_Q^2$ and the affine lines in the $q+1$
directions defined over $\mathbb F_q$, where $Q=q^R$.
\end{enumerate}
Let $o$ be the barycenter of the point sites. Let $A$ and $B$ be
the mean squared distances to $o$ on the point and line sides,
respectively, with uniform weights on each side. Then
\[
                 e\ge (1-q^{-1/2})(A+B).
\]
\end{lemma}

\begin{proof}
The field $\bar k$ contains every $\mathbb F_Q$ in the statement
and satisfies $\omega_{\bar k}\le\omega_k$. To apply
Theorem~\ref{thm:incidence} there, recall that scalar extension
preserves the possible joint spectral evaluations on any finite
placement of $k$-tensors (Section~\ref{sec:prices-background}).
In particular, each finite subset of our metric space has an isometric
realization in a complete spectral space over $\bar k$. Take a pointed
ultralimit of these realizations over the finite subsets containing a
fixed basepoint, using an ultrafilter containing every tail under inclusion.
Each original site is eventually present with all its prescribed
distances. Completeness and uniqueness of geodesics make the resulting
isometric image closed and convex. For a bounded finite distribution
in the ultralimit, the barycentric variance inequality identifies its
barycenter with the limit of the coordinate barycenters. For a
distribution in the image, nearest-point projection places its
barycenter in the image as well. It therefore suffices to prove the
inequalities on placements over $\bar k$.

For the affine plane, apply the unit-mass estimate in the proof of
Corollary~\ref{cor:inc-link} to the configuration $\mathbb F_Q^2\subset\bar k^2$, retaining $q+1$
directions. It bounds the cross term in squared-distance comparison
by $2\sqrt{AB/(q+1)}$, giving the assertion.

For the projective space, map $\bar k^\ell$ linearly into $\bar k^3$,
preserving the ranks of sets of up to three projective
representatives. Such maps form a nonempty open set. Indeed, preserving
each rank is a nonempty open condition and there are only finitely many
conditions. In the resulting projective plane choose an affine chart
containing all the images. Preservation of pairs and triples makes the
images distinct and preserves exactly which triples are collinear.
The secants of the affine configuration thus represent the lines of
the original projective space.

There are $j$ lines through each point and $s=q+1$ points on each line.
Write $r_i=d(o,x_i)$ for point sites and $r_L=d(o,y_L)$ for line sites.
Set
\[
 K_L=\sup_{z\ne o}\Bigl(\sum_{i\in L}
       r_i\cos\angle_o(x_i,z)\Bigr)_+ ,
\]
with zero summands at $r_i=0$. The singleton space satisfies the
claim directly. In every other case Theorem~\ref{thm:incidence}
gives $\sum_L K_L^2\le (j-1)\sum_i r_i^2$. If $M$ is the number of
incidences, triangle comparison and Cauchy--Schwarz bound the mean
cross term by
\[
 \frac{2}{M}\sum_L r_L K_L
 \le 2\sqrt{\frac{j-1}{js}AB}
 \le q^{-1/2}(A+B).
\]
Here $M$ is both $j$ times the number of points and $s$ times the
number of lines. Subtracting the cross term from $A+B$ proves the claim.
\end{proof}

\subsection{Permutations on short address sets}

Choose an integer $\ell\ge3$ and reserve $m$ auxiliary coordinates,
where $m$ will depend only on $q$. For $n=2^h$ put
\[
 r=\left\lfloor \frac{\log_q n}{\ell+m}\right\rfloor,\qquad
 U=\mathbb F_q^r,\quad N=q^r,\quad
 c=\left\lceil n/N^{\ell+m}\right\rceil,\qquad
 \mathcal I=U^\ell\times[c]\times U^m .
\]
Here $[c]=\{1,\ldots,c\}$. Identify the original $n$ addresses with
the first $n$ labels in a fixed ordering of $\mathcal I$ and extend
the input permutation $\pi$ by fixed points. The padded address set
satisfies $n\le|\mathcal I|<2n$.

A \emph{permutation function} takes a height $h$ and a valid input
permutation as parameters and outputs a permutation of $\mathcal I$.
We require a single deterministic program using $nh^{O(1)}$ operations
on words of $O(h)$ bits, with the usual word operations of polynomial
bit complexity. These functions form a countable group $G$ under
pointwise composition. An explicitly listed permutation has a
linear-time inverse, and $\pi$ itself is an element of $G$.

Identify $U^\ell$ with $\mathbb F_q^\ell\otimes U$. For a fixed
subspace $V\le\mathbb F_q^\ell$ of dimension at most two, let $K_V$
be the subgroup preserving every fiber of the projection
$\mathcal I\to(\mathbb F_q^\ell/V)\otimes U$. All auxiliary labels
belong to the fiber. Its size is $cN^{m+\dim V}$, so every hop in
$K_V$ permutes sets of size at most $2^{th+O(1)}$, where
\begin{equation}\label{eq:sort-hop-height}
                         t=\frac{m+2}{m+\ell}.
\end{equation}
The additive constant follows from $c\le q^{m+\ell}$. Increasing
$\ell$ shortens the address within a fiber.

The subgroups $K_V$ generate $G$. We give the coordinate factorization
of a rearrangeable network \cite{benes} in algorithmic form because
the routing protocol will use its factors. Treat $[c]\times U^m$ as
one coordinate. To eliminate a coordinate of size $a$, form a bipartite
multigraph on the remaining input and output names, with one edge per
requested transfer. Each vertex has degree $a$. Give the edges $a$
colors, using each color once at each vertex. In the first stage,
permute the chosen coordinate to its edge color. In the last stage,
permute that color to the destination coordinate. Between these stages
the remaining coordinates are permuted separately for each color.
Iterating gives a bounded number of stages for fixed $\ell,m$.
Each ordinary coordinate stage is in its coordinate line subgroup,
and each auxiliary stage is in $K_0$.

For completeness, the coloring takes deterministic time linear in
the edge count up to logarithmic factors (see also \cite{color}).
An even-regular bipartite multigraph splits into two regular halves
by alternating colors around Euler tours. At an odd degree $a$ we first
remove one perfect matching. To find it, let $v$ be the number of
vertices on each side and take the least $H$ with $2^H>av$. Multiply
each edge multiplicity by $\lfloor 2^H/a\rfloor$ and complete the
graph to degree $2^H$ with copies of an arbitrary dummy matching. The
total dummy multiplicity is less than $av<2^H$. Split into regular
halves repeatedly, keeping the half with fewer dummy copies and
breaking ties deterministically. Each split divides even multiplicities
equally and alternates the remaining odd copies around Euler tours,
distinguishing dummy from original edges. After $H$ splits the retained
degree-one graph has no dummy copy. Only the original edges and the
dummy matching need be stored, with multiplicities fitting in a word.
After the matching is removed the residual even degree splits in half,
so the full coloring recursion has logarithmically many levels.

\subsection{Uniform displacement}

A bound obtained by adding displacements along the coordinate
factorization would depend on $\ell$. The fixed point in the next
lemma bounds the displacement of every product of hops by a single
absolute constant.

\begin{lemma}\label{lem:sort-geometry}
Suppose $G$ acts isometrically on a complete $\operatorname{CAT}(0)$
space $Y$ in which the inequalities of Lemma~\ref{lem:sort-inc}
hold for all placements. There is an absolute constant $C_0$ such
that, if $d(x,gx)\le1$ for all $g\in\bigcup_V K_V$, then $G$ fixes
a point within $C_0/2$ of $x$. In particular $d(x,\pi x)\le C_0$.
\end{lemma}

\begin{proof}
Use a triangle for each flag $0\subset L\subset P$, where $L$ is
a vector line and $P$ a vector plane in $\mathbb F_q^\ell$. We use
the link-energy argument of Garland and its metric fixed-point
version \cite{garland,izekinayatani}. Choose the flag uniformly and,
within it, a vertex or edge uniformly. A placement assigns each $V$
a point $y_V$ fixed by $K_V$. For each $V$ we seek another such point
$z_V$ so that, in a uniform flag through $V$,
\begin{equation}\label{eq:sort-flag-link}
 \mathbb E\,d(z_V,y_W)^2\le \tfrac34\,
                           \mathbb E\,d(y_{W_1},y_{W_2})^2 .
\end{equation}
The vertices $W_1,W_2$ lie opposite $V$, and $W$ is uniform among them.
At $V=0$, the opposite edges are the incidences of projective points
$L$ and projective lines $P$. The barycenter of the $y_L$ is fixed
by $K_0$. Lemma~\ref{lem:sort-inc} bounds its mean squared distance
to a neighbor by $1/[2(1-q^{-1/2})]$ times the opposite-edge energy.
At $V=L$, the barycenter of the $y_P$ with $P\supset L$ is fixed by
$K_L$. Its variance inequality with test point $y_0$ gives the
coefficient $1/2$.

\smallskip\noindent
\emph{The plane vertex.}
When $V=P$, the opposite edge joins $0$ to a line $L\subset P$.
Here $y_0$ and the $y_L$ need not be fixed by $K_P$. We use translates
of these sites to build an affine incidence graph of the same edge
energy. The point barycenter of this graph will be nearly fixed by
$K_P$.

Choose coordinates $(x_1,x_2)\in U^2$ along $P$ in each fiber
modulo $P$, using a fixed basis and linear complement. Each coordinate
pair names a class of $cN^m$ labels. Consider any other partition of
each fiber into classes with these same names and cardinalities,
computed uniformly at all heights and inputs in the time defining
$G$. It is obtained from the coordinate partition by some $g\in K_P$,
by sorting and matching equal names within fibers. Explicitly, if $p$
is the coordinate name, the transported name is $p\circ g^{-1}$.
Assign the new partition the site $g y_0$. Any other such $g$ differs on the right by an element
of $K_0$, so this site is well defined. Projecting the names modulo
$L\otimes U$ gives a coarser partition to which we assign $g y_L$.
The stabilizer of the reference projected partition is $K_L$.
Both assignments commute with
the action of $K_P$. Thus a partition and its projection have site
distance exactly $d(y_0,y_L)$.

Let $Q=q^R$, where $R\le m$ will be chosen below. Use $R$ auxiliary
$U$-coordinates as the values of an $\mathbb F_q$-linear map
$u:\mathbb F_Q\to U$ on a fixed basis of $\mathbb F_Q$. For a uniform
label in a fiber, $u$ is uniform independently of $(x_1,x_2)$. For
each $(a,b)\in\mathbb F_Q^2$ form the partition with name
$(x_1+u(a),x_2+u(b))$. The name is uniform on $U^2$. Along an affine
line of direction $L$ defined over $\mathbb F_q$, the name modulo
$L\otimes U$ is constant as a function of $(a,b)$, by linearity of
$u$. We therefore obtain sites for all points and all retained lines.
The incidence energy is $e_P=\mathbb E_{L\subset P}d(y_0,y_L)^2$.
At the point barycenter $b_0$, Lemma~\ref{lem:sort-inc} now gives
$(1-q^{-1/2})(A+B)\le e_P$.

\smallskip\noindent
\emph{A fixed center.}
For $d_*$ independent samples from the point-site distribution, form
a successive geodesic mean by moving a fraction $1/i$ toward sample
$i$ at step $i$. Its final value $M_{d_*}$ satisfies
\begin{equation}\label{eq:sort-mean}
                  \mathbb E\,d(M_{d_*},b_0)^2\le A/d_*.
\end{equation}
Indeed, condition on the previous mean. Squared-distance convexity
and the variance inequality at $b_0$ bound the new expected squared
error by $(1-1/i)^2$ times the old error plus $A/i^2$.

Call a tuple of partitions independent if its joint names are uniform
in every fiber at every height and input. This is a property of a
realized tuple. Choose a large absolute integer $d_*$ and set $m=R=4d_*$.
For $4d_*$ independent uniform samples from $\mathbb F_Q$, the probability
of a linear dependence over $\mathbb F_q$ is at most
$\sum_{j=0}^{R-1}q^{j-R}<1/(q-1)$. With probability at least $1/2$, the
two coordinates of each of $2d_*$ sampled affine points give $4d_*$
linearly independent field elements. Evaluation of $u$ on them is
then jointly uniform even conditional on $(x_1,x_2)$, so the $2d_*$
partitions are independent.

The group $K_P$ is transitive on the efficiently specified independent
tuples of this length, by matching joint names within each fiber.
Consequently the distance between the means of the two halves is the
same for all such tuples. Its unconditional expected square is at most
$4A/d_*$ by \eqref{eq:sort-mean} and the triangle inequality. For
each independent tuple the square is therefore at most $8A/d_*$.
The same conditioning gives an independent $d_*$-tuple of affine
partitions whose mean $M$ satisfies $d(M,b_0)^2\le2A/d_*$.

Fix $g\in K_P$. The tuple giving $M$ and its $g$-translate may be
dependent. Construct a third tuple jointly independent with each of
them separately. In each fiber join the old name at a label to its
translated name by one edge for that label. Both sides of the resulting
bipartite graph have $N^{2d_*}$ vertices, each of degree
$cN^{m+2-2d_*}$. The degree is divisible by $N^{2d_*}$. Color its
edges properly as above and group the colors into $N^{2d_*}$ sets of
equal size. Regard the color set of an edge as its third name in
$U^{2d_*}$. At every vertex, each third name occurs equally often,
which proves the two required independence statements. The coloring
also supplies the third tuple in the time defining $G$.

The mean of this tuple lies within $\sqrt{8A/d_*}$ of both $M$ and
$gM$. Hence
\[
 d(b_0,gb_0)\le 2\bigl(\sqrt{2A/d_*}+\sqrt{8A/d_*}\bigr)
                      \le 0.04\sqrt A
\]
on choosing $d_*$ sufficiently large. The circumcenter $z_P$ of the
$K_P$-orbit of $b_0$ is fixed by $K_P$ and lies within $0.04\sqrt A$
of $b_0$. With half the mass on each side of the affine incidence graph,
the root mean squared distance to $z_P$ is at most
$\sqrt{(A+B)/2}+0.04\sqrt A$. In this average an affine point site is a
translate of $y_0$, and an affine line site of direction $L$ is a
translate of $y_L$. Since $z_P$ is fixed, these distances are precisely
the distances from $z_P$ to $y_0$ and $y_L$. The inequality
$(1/\sqrt2+0.04)^2/(1-q^{-1/2})<3/4$ proves
\eqref{eq:sort-flag-link} at $P$.

\smallskip\noindent
\emph{Iteration on the flags.}
Initialize $y_V$ at the circumcenter of $K_Vx$, within distance one of
$x$. Replace each $y_V$ by its midpoint with $z_V$, which is still fixed
by $K_V$. Write $e$ for the mean squared edge length of a flag. By
convexity of distance between geodesics, the squared distance between
the new endpoints at $V,W$ is at most
$\bigl(d(y_V,z_W)^2+d(z_V,y_W)^2\bigr)/2$.
Averaging \eqref{eq:sort-flag-link} over flags thus gives
$e_{\rm new}\le3e/4$. By the triangle inequality through a random old
neighbor, the root mean squared movement per vertex is at most
$(\sqrt e+\sqrt{3e/4})/2\le\sqrt e$.
Initially $e\le4$. The movements are summable in the finite weighted
metric product, and each edge in its limit has length zero. The flag
graph is connected, so the limiting placement has a single value
$y_*$. This value is fixed by all the $K_V$ and hence by $G$.
Summing the movements gives $d(x,y_*)\le1+2/(1-\sqrt{3/4})$.
We may take $C_0=100$.
\end{proof}

\subsection{Finite tests on permutation orbits}

We now return to the finite set of addresses. In the routing protocol
a proposed price is assigned separately to each payload wire. We will
test these prices in balls of bounded radius. To compare price
realizations on overlapping balls we must also test the other tensors
on their common sites.

Take a finite symmetric set of elements of $\bigcup_V K_V$ containing
the factors of a fixed coordinate factorization of $\pi$ and their
inverses. Call these the \emph{core generators}. For each input they
define an undirected graph on $\mathcal I$, with a \emph{core link}
for the hop from a label to its image under a generator. Its degree
is bounded by the size of the set. Each $i$ has a path to $\pi(i)$
along the factors, of length bounded by a fixed rational $a_{\mathrm p}$
independently of $h$ and the input.

The four test parameters are the core generators, a ball radius,
a maximum number of sites, and a maximum local dimension. Give each
payload wire its own commodity index. Also index every nonzero $k$-tensor
of permitted format on an ordered list of distinct sites in a common
test ball, using fixed coordinate bases. Assign this index on the
whole graph, so two balls containing the same tensor on the same
ordered sites see one price. There are $O(|\mathcal I|)$ indices
because each ball has bounded size and the field is finite.

Assign each commodity a nonnegative price. Cap a payload price at
$a_{\mathrm p}$. For each other format choose a uniform rational cap
sufficient to create the tensor at one site and transport its legs in
binary coordinates along core links in the ball. On a test ball compare
the assigned prices of all its commodities with an ordinary spectral
price on its sites. The error is the largest absolute discrepancy on a
commodity or positive excess above unit cost on a core link with both
endpoints in the ball. Duplicate commodities are compared separately
with the spectral price of the same tensor. The minimum error over ordinary prices, allowing the collapsed price,
is the \emph{discrepancy} of the test. This minimum is attained.
Indeed, the core links connect the ball, so a spanning tree bounds
the scale of any price with bounded error, and prices at bounded scale
form a compact set. On a singleton use the collapsed price. A center
\emph{passes} to tolerance $\rho$ if this minimum is at most $\rho$.

\begin{lemma}[Finite detection]\label{lem:sort-local-tests}
For each $\ell$ there are finite test parameters and $\rho>0$ such
that the following holds for all sufficiently large $h$ and every
valid permutation input. If all but a fraction $\rho$ of the centers
pass to tolerance $\rho$, the assigned payload prices satisfy
\[
        |\mathcal I|^{-1}\sum_{i\in\mathcal I} x_{i,\pi(i)}
                  \le C_0+1 .
\]
Here $C_0$ is the absolute constant of Lemma~\ref{lem:sort-geometry}.
\end{lemma}

\begin{proof}
If the assertion fails, choose a sequence of counterexamples with
$h\to\infty$. Enlarge the finite symmetric sets of core generators
to exhaust the countable union of the $K_V$, and let the radii and format
bounds tend to infinity. Take the tolerance and permitted exceptional
fraction tending to zero, keeping the original factorization of $\pi$
and the payload cap fixed. We use these counterexamples to construct
an action in which each hop has displacement at most one and the
target permutation has displacement at least $C_0+1$.

\smallskip\noindent
\emph{The marked sites.}
At each passing center $i$ take a complete $\operatorname{CAT}(0)$
spectral realization of a price attaining the test. At every other
center use the one-point space. For $g\in G$ mark the site $g^{-1}i$
in the chart at $i$. A fixed word in subgroup elements represents
$g^{-1}$, so eventually these hops are tested and the ball contains
the whole path from $i$ to $g^{-1}i$. Thus the distance from the
center to this mark is uniformly bounded for each fixed $g$.
Use the center itself when the mark is unavailable.

At each index take the metric product over all centers with the
root mean square distance, pointed at the vector of centers.
Take a pointed ultralimit of the products and let $Y$ be the closed
convex hull of the resulting marks. The product and ultralimit are
complete $\operatorname{CAT}(0)$ spaces, as one sees from the
midpoint inequality. Both incidence bounds of Lemma~\ref{lem:sort-inc}
hold in $Y$. A product barycenter is coordinatewise, so the squared
inequalities average over the coordinates. In the ultralimit the
variance inequality identifies the limit of barycenters with the
barycenter of a bounded finite placement. Finally, projection onto a
closed convex subset preserves the barycenter of a distribution in
that subset.

\smallskip\noindent
\emph{Translation of the marks.}
We claim that left multiplication $g\mapsto ug$ extends to an
isometry of $Y$. A point of a dense subset of $Y$ is given by a
finite expression of geodesic interpolations of marks. Fix two such
expressions. At center $i$ the translated expressions have the
same underlying address labels as the original expressions at center
$u^{-1}i$, since $(ug)^{-1}i=g^{-1}u^{-1}i$.

For sufficiently large indices, both balls contain these labels.
Both centers pass outside a fraction tending to zero, since
$i\mapsto u^{-1}i$ preserves uniform measure. For every fixed tensor
on the distinct labels the two realizations then differ by at most
twice the tolerance once the format bounds include that tensor.
Their evaluations use its common commodity index. It follows that
the distances between the two expressions in the two charts differ
by a quantity tending uniformly to zero on such pairs. Indeed, a
violating sequence has a subsequence with a fixed pattern of coincident
labels. The finitely many factor words bound the diameter of the
marked list in both charts, so the tree estimate bounds their price
scales. By compactness take limiting prices on the list. They agree
on every tensor by the format exhaustion. Continuity of distances
between geodesic expressions, from Lemma~\ref{lem:geodesic}, gives
equal limiting distances, a contradiction.

All expression distances are bounded even on exceptional pairs.
In a passing chart the fixed word lengths bound all the marks, and
a failed chart has collapsed to a point. The exceptional pairs
therefore make a vanishing contribution to the squared mean. Reindexing
by $i\mapsto u^{-1}i$ proves the claimed equality of distances in the
limit. The assignment on expressions thus extends to $Y$. Its inverse
comes from $u^{-1}$, and the group law holds first on expressions and
then by continuity on $Y$.

Let $x$ be the identity mark. For each individual $g\in K_V$ its hop
and inverse eventually belong to the tested set, so $d(x,gx)\le1$.
The mark $\pi^{-1}x$ at center $i$ corresponds to the address $\pi(i)$.
In a passing chart its distance from the center agrees up to the
tolerance with the assigned price of the payload wire $i\to\pi(i)$.
In a failed chart the distance is zero while the price is at most
$a_{\mathrm p}$. Thus $d(x,\pi^{-1}x)$ is the ultralimit of the root
mean square payload price. This is at least the limiting lower bound
$C_0+1$ on the arithmetic means of the counterexamples.
Since $d(x,\pi^{-1}x)=d(x,\pi x)$, this contradicts
Lemma~\ref{lem:sort-geometry}.
\end{proof}


\section{Bounded-width permutation routing}\label{sec:sort-route}
\subsection{From a failed price test to a routing identity}

Fix the core generators and the test parameters of
Lemma~\ref{lem:sort-local-tests}. Write $M=|\mathcal I|$, so that
$n\le M<2n$. Sending a bit along the designated path from $i$ to
$\pi(i)$ pays for every hop on that path. We seek an exact sum of
networks in which the average transport cost per payload bit has an
absolute bound, even as the parameter $\ell$ increases.

\begin{proposition}\label{prop:sort-protocol}
Suppose $\omega_k<9/4$ and fix $H\ge1$. For all sufficiently large
$h$, every valid set of tags and every $h\le w\le h^H$, the associated
permutation tensor for $w$-bit payloads is an exact sum over $k$ of
networks with the following properties.
\begin{enumerate}
\item Each label has a bounded number of fixed slots per core generator.
A used slot carries $O(w)$ bits to a neighbor along the generator or its
inverse. Both ends belong to the same prescribed group of at most
$2^{th+O(1)}$ labels. The sum of the carried widths in any summand is
at most $Cnw$, where $C$ is independent of $\ell$ and $t$.
\item At each label a fixed trivalent network of size
$O(w(1+\log w)^{O(1)})$ implements the local maps. Its configuration
consists of a bounded list of menu indices, widths and offsets.
\item A summand is specified by $O(n\log(2w))$ bits. Given these bits
and the tags, a deterministic controller computes the local
configurations and matchings in $nh^{O(1)}$ word operations.
The exponents and slot capacities are independent of $H$. The running
time may have a multiplicative constant depending on $H$.
\end{enumerate}
\end{proposition}

\begin{proof}
We use the local replacement argument of
\cite[Section~\ext{G}{sec:protocol}]{oaG}. A \emph{request} for a
commodity leaves that tensor to be supplied at specified legs.
Its cap $a_i$ bounds the cost of supplying it directly, by preparing
the tensor at one site and sending its legs along designated core
paths. We encode each leg in binary and pad its unused coordinates
by zero. Multilinearity permits us to substitute a tensor identity
for several requests at any stage of a contraction.

\smallskip\noindent
\emph{The identity on a ball.}
Let $x\in\prod_i[0,a_i]$ be a price array with positive discrepancy on
a ball. Request one copy of each commodity with $x_i>0$ on the ball.
Their tensor product $D$ has assigned total price $L=\sum_{i:x_i>0}x_i>0$,
since the zero array passes the test. The available supplies $T_j$
are all the commodities on the ball, at their assigned prices, and
its binary core links, at price one. On a singleton we can prepare
every demand locally. On a larger ball there are prices $l_j$ strictly
above the supply prices and a number $0<L_-<L$ such that every
normalized spectral log price $P$ satisfies
\[
                 P(D)<L_-\max_j\frac{P(T_j)}{l_j}.
\]
Indeed a spanning tree of binary core links supplies the full binary
diagonal, of normalized log price one. Thus the maximum ratio is at
least the reciprocal of the sum of the $l_j$ over the tree. If the
strict inequality with $L$ failed for perturbations decreasing to
the original supply prices, rescale each violating $P$ by dividing
by its maximum ratio. The tree bound keeps the scales bounded.
A limit price then satisfies every supply bound and gives $D$ price
at least $L$. Additivity on products forces equality at each
positive-priced commodity, and the remaining commodities have price
zero. This limit would pass the test on the ball. For a fixed
successful perturbation, compactness permits the decrease from $L$
to $L_-$ with a uniform strict margin.

Choose an integer $b_0$ and form the direct sum
\[
               U=\bigoplus_j T_j^{\otimes\lfloor b_0L_-/l_j\rfloor}.
\]
Rounding costs a bounded amount in log price. For sufficiently large
$b_0$, every normalized spectral price gives $U$ a higher value than
$D^{\otimes b_0}$. Strict spectral comparison in
Section~\ref{sec:prices-background} now gives an exact restriction of
$D^{\otimes b_0r}$ from at most $K(b_0)$ direct copies of
$U^{\otimes r}$ for every $r\ge1$. Expand the direct sums and apply
the restricting maps. The resulting ordinary sum supplies $b=b_0r$
copies of every demand. Each term uses supplies of total price at
most $bL_-$. Its index specifies a direct copy and an ordered word
of $r$ supply choices. With $s$ supply types this takes at most
$\lceil\log_2(K(b_0)s^r)\rceil$ bits. The ratio of this length to
$b$ can be made arbitrarily small by first increasing $b_0$ and then
$r$.

Repeat the identity $J$ times. Use only the \emph{histogram} of the
$J$ term indices as the outer selector. This costs $O(\log(2J))$
bits because the alphabet is fixed. Within the corresponding network
a root sums over ordered words, sends each word along the core tree
and enforces its prescribed histogram. The histogram fixes the
number of copies of every supply. The word assigns these copies to
the individual applications of the restricting maps. Each ordered
word contributes exactly once with its original coefficient over
$k$. Choose the rate of the term indices low enough that the tree
transmission costs at most $Jb(L-L_-)/2$. The replacement thus saves
a positive amount per block of demands, after paying for its own
control bits.

These demands remain positively priced, and the replacement retains
a strict saving, throughout a neighborhood of the original array.
There are finitely many ball types since the field, radius, degrees
and formats are fixed. Compactness of the price arrays of discrepancy
at least $\rho/2$ gives a finite menu of replacements. On this set,
the saving per unit of demand and the selected demand prices both
have uniform positive lower bounds.

\smallskip\noindent
\emph{Bounding the total width.}
Follow one summand in the expansion. Let $wz_i$ be its integer stock
of outstanding requests for commodity $i$, and let $S$ be its
installed core width divided by $Mw$. Give the summand the account
$S+F(z)$, where
\[
       F(z)=\frac1M\max_{0\le x\le a}
                   \sum_i\bigl(z_ix_i-\tfrac12\mu x_i^2\bigr).
\]
There are $O(M)$ coordinates with uniformly bounded caps. Choose
a positive rational $\mu$ such that $\mu\sum_i a_i^2/(2M)\le1$.
The maximizing price $x_i=\min(a_i,z_i/\mu)$ satisfies
$z_i\ge\mu x_i$. Its Lipschitz bound gives, for $z+v\ge0$,
\[
       F(z+v)-F(z)
       \le \frac1M\sum_i x_iv_i+\frac1{2\mu M}\sum_i v_i^2 .
\]
Fix an integer $B>\max_i\mu a_i$. Whenever a stock exceeds $wB$,
supply the excess directly. On the interval $z_i\ge B$ the
marginal price equals $a_i$, so this operation does not increase
the account.

If the mean maximizing price of the $M$ payloads exceeds $C_0+2$,
Lemma~\ref{lem:sort-local-tests} gives more than $\rho M$ balls of
discrepancy greater than $\rho$. A greedy packing selects $\Omega(M)$
disjoint such balls, allowing a fixed slack in the test. Indeed the
discrepancy is $1$-Lipschitz in the array of prices. A sufficiently
fine rational grid can assign a valid menu entry to each cell meeting
the set of discrepancy at least $\rho$, with the entry valid throughout
the cell. These finite tables suffice to select and pack eligible balls.

On each packed ball apply its replacement with
$J=\lfloor\alpha w/b\rfloor$, for a small fixed rational $\alpha>0$.
The bound $z_i\ge\mu x_i$ and the positive lower bound on selected
demand prices ensure sufficient stock. Charge each new core link to
$S$, and enter other supplies as new requests. In the linear bound on
the account, each ball saves a fixed positive amount per retired demand
block, including transmission of the word. For large $w$ the packed
balls therefore save at least $c\alpha$ in the normalized account for
some $c>0$. The finite menu and the disjoint balls give
$\max_i|v_i|=O(\alpha)$, so the quadratic error is $O(\alpha^2)$.
Fix $\alpha$ small enough that each parallel round lowers the account
by a fixed positive amount in every resulting summand. Supply any
excess over the stock cap after the round.

Starting from zero, introduce all $w$ copies of each payload in equal
increments on a fixed fine mesh, rounding to integers with exact total
$w$. Before each increment use parallel rounds until the mean payload
price is at most $C_0+2$. The account is nonnegative. The increments
have total linear cost at most $C_0+2$, and their summed quadratic
cost is at most one for a sufficiently fine mesh and sufficiently
large $w$. Supply excess stock after each increment. Every summand
therefore undergoes only boundedly many rounds, and its final account
is at most $C_0+3$. Supply all the remaining requests directly.
Taking $x=a$ in the definition of $F$ gives
\[
             \frac1M\sum_i a_i z_i\le F(z)+1.
\]
The total width is at most $(C_0+4)Mw\le 2(C_0+4)nw$. This last
coefficient is absolute.

\smallskip\noindent
\emph{Realizing the summands.}
A label lies in at most one packed ball per round. Every direct
supply that uses a given edge originates within bounded distance of
it. Bounded degree, boundedly many commodity types per neighborhood,
and the stock cap thus give boundedly many slots of width $O(w)$
per generator. The generator prescribes the group containing both
ends of the slot.

At the sites of a commodity, keep its requests in common order,
retiring them from the front and appending new requests at the back.
For each replacement the histogram determines the lengths of the
lists. Inside its network, prefix counts of the ordered word assign
supply positions to repeated applications. Comparison sorting on
addresses of $O(1+\log w)$ bits realizes these permutations and
partial matchings with $O(w(1+\log w)^{O(1)})$ trivalent gates at a
site. Complete partial matchings deterministically and set idle
coordinates to zero. Each individual block map has fixed size over
$k$ and admits a constant-size trivalent implementation. Only
boundedly many local fields select the maps, lengths and offsets.

The controller enumerates the bounded neighborhoods, computes prices
from the stocks, uses the finite tables to pack eligible balls, and
assigns the slots. There are $O(n)$ histograms over the boundedly
many rounds, with a unique encoding of $O(\log(2w))$ bits each.
Enforce valid encodings and zeros in unused selector positions.
Only local counts and indices enter the controller, while the ordered
words are summed inside the networks. The enumeration and stock
updates take $nh^{O(1)}$ word operations on tags of $h$ bits and
counts of $O(\log(2w))$ bits. This proves the three bounds.
\end{proof}

\subsection{From local transports to stable sorting}
\label{subsec:sort-recursion}

Write $S(h,w)$ for the least trivalent size of $\mathrm{Tag}(h,w)$
over $k$, and put $n=2^h$. Proposition~\ref{prop:sort-protocol} bounds
the transport used in each summand. To derive a recurrence for $S$,
we must fit the calls from all summands into a single network.
Different summands can use the same call if their data can reach it
at low cost.

\paragraph{Sharing a fixed pool of calls.}
Fix the dimension $\ell$, the address reduction factor
$t=(m+2)/(m+\ell)$ and the protocol of Proposition~\ref{prop:sort-protocol}.
In each term the protocol has at most $Ln$ links of total width at most
$Cnw$. Of these two constants, $C$ is independent of $\ell$.
Choose $0<Y\le1$ with $L\le C/Y$ and split each link into chunks of
$w'=\lceil Yw\rceil$ bits. Pin the padding of a final chunk to zero
at both ends. The number of nonempty chunks in any term is at most
\[
             Cnw/w'+Ln\ \le\ m_*n,\qquad m_*=2C/Y .
\]
For this fixed protocol there are $O(n)$ candidate slots, grouped
according to the short-address sets within which they can communicate.
Duplicating a group for the boundedly many link ends or chunks it
handles changes its size only by a constant factor.

We use the assignment graphs of
\cite[Lemma~\ext{G}{lem:slack-compaction}]{oaG}. Such a graph connects
at most $cK$ candidates to $K$ slots. Every active subset of at most
$K/2$ candidates can be assigned injectively to adjacent slots.
The degrees are bounded in terms of $c$, and the assignment takes
$O_c(K)$ deterministic word operations with the graph as fixed advice.

Choose a dyadic capacity $U=2^{th+O(1)}$ at least twice the size of
the largest candidate group. Pack whole groups greedily into lists of
capacity $U$, closing a list when the next group would exceed that
capacity. Each closed list is at least half full, so there are
$O(n/U+1)$ lists. Fix a small dyadic fraction $K/U$. In a list,
divide its active chunks in order into batches of size $K/2$, allowing
a shorter last batch. With $K$ slots reserved for each nonempty batch, the total reservation
is at most
\[
             2m_*n+O((n/U+1)K)\ \le\ 3m_*n .
\]
Here $U=o(n)$ and we choose $K/U$ after fixing the packing constants.

For each list install enough candidate batches to handle every chunk
in it. Use one assignment graph to place the source ends of a batch in
$K$ slots and another for the target ends. A second assignment step
handles whole batches. Its targets are a shared pool of
$\lceil6m_*n/K\rceil$ physical modules, each with $K$ input and output
positions. At most half the pool is active, and the number of candidate
batches is a bounded multiple of the pool size. Thus the two lists of
each active batch can be moved together to its assigned module.
Including idle modules, at most $7m_*n$ positions are installed for
all large $h$.

Every module uses $\mathrm{Tag}(h',w')$ with $h'=\log_2 K=th+O(1)$.
Fill unmatched positions in order, pin unused bits to zero and give
an idle module the identity tags and zero payloads. The assignments
and padding are deterministic once the term and its tags are given.
Because the assignment graphs have bounded degree, the switches cost
$O(nw)$ in total, including those that move two lists of $Kw'$ bits
at a time.

The remaining control uses only tags, selectors and slot descriptions.
It can be checked by guessing a time-stamped execution, verifying each
instruction, and comparison-sorting all memory transactions by address,
time and access order within each instruction to check every read against the last write. Fixed initial
memory and padding make the accepting execution unique for every legal
selector. A conventional comparison network also checks the original
tags for distinctness and rejects invalid permutations. For each fixed
$H$ and $h\le w\le h^H$, the size outside the shorter calls is
\begin{equation}
  D_{\rm control}(h,w)=
  O_H\bigl(nw(\log(2h))^{O(1)}+n h^J\bigr),\qquad J>2 .
  \label{eq:sort-control}
\end{equation}
The first term pays for the local maps with their polynomial overhead
in $\log(2w)$. The second pays for $nh^{O(1)}$ deterministic word
operations, their bit implementation and the memory check. The
exponents can be chosen independently of $H$. We have obtained the
recurrence
\[
 S(h,w)\ \le\ D_{\rm control}(h,w)
       +\frac{7m_*n}{2^{h'}}\,S(h',w').
\]

\paragraph{Wide payloads.}
It suffices to prove the estimate for $0<\zeta<1$.
Choose $\ell$ large enough that $14C t^\zeta<1/4$, and then fix the
resulting $Y$ and $J$. Shortening an address from $h$ to $h'$ reduces
the allowed overhead $h^\zeta$ by a factor approaching $t^\zeta$.
To account for calls close to the width at which conventional sorting
takes over, we let the overhead exponent decrease gradually with the
width.

Choose a smooth nonincreasing $u:[1,\infty)\to[\zeta,2]$, equal to
$2$ on $[1,J]$ and eventually equal to $\zeta$. Take such a long
transition interval that $b=-u'$ satisfies $0\le b<1$ and
$Y^{-b}<3/2$ everywhere. Fix $H$ in the interior of the constant
tail. For $h\ge2$ and $h\le w\le h^H$ we claim
\begin{equation}
       S(h,w)\ \le\ D n w h^{u(a)},\qquad a=\log_h w .
       \label{eq:sort-width}
\end{equation}
If $a\le J$, Batcher's bitonic sorting network~\cite{batcher} uses
$O(nh^2)$ comparators of bit cost $O(h+w)$ to put the tagged sources
in destination order. Comparing their tags with the prescribed
positions and their payloads with the proposed outputs gives
\eqref{eq:sort-width}.

For $a>J$ induct on $h$. For large heights $h'<h$ and $w'\ge h'$.
Writing $a'=\log_{h'}w'$, we have uniformly on $J\le a\le H$,
\[
 a'-a=\frac{\log Y-a\log t+o(1)}{\log h},
 \qquad
 \frac{w'(h')^{u(a')}}{w h^{u(a)}}
   =(1+o(1))\,t^{u(a)+a b(a)}Y^{1-b(a)}.
\]
The second formula is Taylor expansion on a compact interval.
If $w'>(h')^H$, split the payload evenly into
$\lceil w'/\lfloor(h')^H\rfloor\rceil$ slices, using the same tags
for each. The number of slices is bounded, their widths sum to $w'$,
and their width exponents are $H+o(1)$, where $u=\zeta$. Copying
the tags costs $O(h')$ per slot and slice, within
\eqref{eq:sort-control}. The same per-slot estimate therefore applies.

After multiplying by the ratio $7m_*=14C/Y$ of installed to original
positions, the recursive calls use at most the fraction
$(1+o(1))14C t^\zeta Y^{-b(a)}<1/2$ of the claimed bound.
The extra cost \eqref{eq:sort-control} is $o(nwh^{u(a)})$ uniformly,
since $a\ge J$ and $u\ge\zeta$. Enlarging $D$ to cover bounded
heights closes the induction. For $w>h^H$, split into slices as
evenly as possible with maximum width $\lfloor h^H\rfloor$. Each
width exponent is $H+o(1)$ for large $h$, and tag copying costs at
most a constant times the total data volume. At bounded heights the
bitonic network costs $O_h(w)$. This proves, for some $A>1$,
\begin{equation}
         S(h,w)=O_\zeta(2^h w h^\zeta)\qquad(w\ge h^A).
         \label{eq:sort-all-wide}
\end{equation}

\paragraph{From tagged permutations to stable sorting.}
For $2^s$ records of $v$ bits, a stable sort uses one
$\mathrm{Tag}(s,v)$ and $O(2^s(s+v))$ additional gates. Give each
proposed output a guessed destination in the original input list.
Use the tagged permutation to check that this sends every output back
to an equal input. Along the output list check strictly increasing
pairs (key, guessed input position) in lexicographic order. All records
with a common key must then be exactly the corresponding input
subsequence. This forces the order and every guessed index, even for
identical records. The stable-sort relation thus inherits the
wide-record bound.

\paragraph{Packing ordinary-width records.}
We use the wide-record bound twice in each radix pass.
First take $\zeta=1/2$ in \eqref{eq:sort-all-wide}, with threshold
$a>1$. Then take $\zeta=1/(2a)$ there, with threshold $B>1$.
Start with $O(n)$ records whose keys and payloads have $O(h)$ bits.
Choose a radix digit of
$s=\lfloor h^{1/a}\rfloor$ bits. Run $O(h/s)$ stable passes starting
with the least significant digit, zero-padding the last digit.
Address lengths for the two sorts will be about $2s$ and $h$
respectively.

For one pass, pad the list to blocks of $2^{2s}$ slots. Sort each
block by a dummy flag and the digit, with dummies last in their old order. Choose a power of two
$P\in[h^{B-1},2h^{B-1})$. Extend each run of a single real digit to
a multiple of $P$, inserting dummies at its end. Each pack of $P$
consecutive slots now contains records of at most one real digit.
When the packs are sorted stably by digit, equal-digit records retain
their order both within a block and across blocks. Empty packs sort
last. Figure~\ref{fig:sort-packs} shows one pass with two blocks.

\begin{figure}[!hbp]
\centering
\begin{tabular}{@{}l@{\qquad}c@{\qquad}c@{}}
\toprule
Input block in order & Packs of digit $0$ & Packs of digit $1$\\
\midrule
$a_1\, b_0\, c_1\, d_1$ & $[b,\cdot]$ & $[a,c]\ [d,\cdot]$\\
$e_0\, f_1\, g_0\, h_0$ & $[e,g]\ [h,\cdot]$ & $[f,\cdot]$\\
\bottomrule
\end{tabular}
\caption{One radix pass with two records per pack. A subscript is the
digit being sorted and a dot is a dummy slot. Stable sorting of whole
packs collects the middle column, then the right column, each from
top to bottom. Reading the real records gives $b,e,g,h,a,c,d,f$.}
\label{fig:sort-packs}
\end{figure}

Within a block a fixed allowance of $O((2^s+1)P)$ extra slots covers
the padding for all real runs, retains the old dummies and rounds the
total length to a multiple of $P$. Scans find the length of each run,
round it up, and assign destinations to all old records, including the
old dummies. A second local tagged permutation checks these destinations.
Append new dummies and require their guessed destinations to be strictly
increasing. Only the unused positions are available to them, so there
is exactly one valid filling. Additional temporary padding to a power of
two is removed after the call. Scan each pack to obtain its digit before
the global stable sort. Dummies, including any inside a pack, are ignored
when reading the real order.

The list remains of size $O(n)$ across all passes. In one pass its
multiplicative growth is $1+O(2^{-s}P)$, with an additive cost
$O(2^{2s}+2^sP)$ for the last block. Over $O(h/s)$ passes the
multiplicative factors have bounded product because $P$ is polynomial
in $h$ and $s=\lfloor h^{1/a}\rfloor$. The accumulated additive cost
is $o(2^h)$ since $s=o(h)$.

Each local call has address length $2s+O(1)$. Padding its records to
a sufficiently large constant times $h$ meets the first width threshold,
so it costs $O(hs^{1/2})$ per slot. The global sort has $O(n/P)$ packs
of $O(hP)$ bits and address length $h+O(\log h)$. Since $hP\ge h^B$,
constant-factor padding meets the second threshold and the global
cost is $O(nh^{1+1/(2a)})$. The scans, copies and order checks
fit the same bound. All passes together cost
\[
            O\bigl((h/s)\,nh^{1+1/(2a)}\bigr)
               =O\bigl(nh^{\,2-1/(2a)}\bigr).
\]

To realize $\mathrm{Tag}(h,h)$, take the union of the $n$ source records
and the $n$ proposed target records. Give each source its claimed
destination and each target its prescribed position, and radix-sort the
union by destination, breaking ties by source before target. Scan the
real records, skipping dummies. For each destination in order the scan
requires one source followed by one target with equal payloads.
An $O(h)$-bit register suffices. This scan rejects repeated or missing
source tags and requires no rearrangement of the final dummy slots.

Thus $\mathrm{Tag}(h,h)$ has exponent $p=2-1/(2a)$, with $1<p<2$.
For width $ch$, split the payload across boundedly many tagged calls
with the same tags, padding when necessary. The reduction from stable
sorting to tagged permutations now yields the ordinary-width assertion
and completes the proof of Theorem~\ref{thm:sort}.
